\chapter{Terminology}\label{ch:terminology}
\begin{description}[style=nextline,leftmargin=0pt]
  \item[Alarm] An operator notification requiring a defined response within a defined time to prevent or mitigate an undesirable consequence.
  \item[Warning] An advisory condition that increases awareness but does not meet the alarm definition.
  \item[Fault] A detected abnormal condition of a device, signal, communication path, or control function.
  \item[Control narrative] A written description of how a control system is intended to operate, including operating modes, sequences, permissives, interlocks, alarms, operator actions, failure responses, and interfaces with other equipment.
  \item[Table of Interlocks and Permissives] A project document that lists the conditions that allow equipment to operate and the conditions that stop or prevent operation. It normally identifies the initiating condition, affected equipment or action, required response, reset behavior, and related notes. The project team may begin the table earlier, but it cannot be finalized until the PHA is complete and its controls-related recommendations have been incorporated.
  \item[Process Hazard Analysis (PHA)] A structured review used by the project team to identify process hazards, evaluate existing safeguards, and recommend additional protection where needed. PHA results may create or change alarms, permissives, interlocks, and trips.
  \item[Permissive] A condition that must be satisfied before an action may start or continue. A permissive is not automatically a safety function.
  \item[Control interlock] Logic that prevents or stops an action to protect normal operation or equipment. It is implemented in the basic process control system unless the risk assessment requires independence.\footnote{EFC currently uses interlocks in the standard PLC to serve as its safety instrumented functions (SIFs); we do not yet design separate, integrity-rated SIFs. This is expected to change when safety PLCs are added to our systems, and this guide will be updated then.}
  \item[Equipment protective interlock] A protective action intended principally to prevent equipment damage. Its required independence and integrity must be established by risk assessment.
  \item[Emergency-stop function] A complementary protective measure intended to avert or reduce an existing hazard. It is not a substitute for safeguarding or risk reduction by design.
  \item[Trip] An automatic transition to a defined safe or protective state.
  \item[Bypass/Override] An authorized temporary defeat or substitution of a protective condition or output. It must be controlled, visible, time limited, and logged.
  \item[Inhibit] Prevention of an alarm or protective action from operating while its initiating condition may still exist.
  \item[Suppression] Prevention of an alarm based on a designed operating state or condition.
  \item[Shelving] A time-limited operator action that temporarily removes an alarm from normal presentation while retaining its record.
  \item[Safe state] The process or equipment condition defined by the risk assessment as producing an acceptable level of risk.
  \item[First-out] Identification and retention of the initiating condition in a group of protective trips.
  \item[Management of Change (MOC)] The controlled review, authorization, implementation, validation, and documentation of a change.
  \item[Short-Circuit Current Rating (SCCR)] The maximum prospective short-circuit current at a stated voltage that an assembly or component can safely withstand under specified conditions. The panel SCCR must be at least equal to the available fault current at its installation point.
  \item[FAT/SAT] Factory Acceptance Test and Site Acceptance Test, respectively.
\end{description}
